\documentclass[12pt]{article}

\usepackage{amsmath}
\usepackage[margin=1in]{geometry}
\usepackage{fancyhdr}
\usepackage{amsthm}
\newtheorem{lemma}{Lemma}
\usepackage{braket}
\usepackage{amssymb}
\usepackage{pgfplots}
\usepackage{tikz}
\usepackage{bbm}
\usepackage{bm}
\usetikzlibrary{shapes.geometric}
\usepackage{enumerate}
\usepackage[shortlabels]{enumitem}

\newcommand{\pderiv}[2][]{\frac{\partial#1}{\partial#2}}

\pagestyle{fancy}
\fancyhead[l]{Hudson Benites}
\fancyhead[c]{Day 7}
\fancyfoot[c]{\thepage}
\renewcommand{\headrulewidth}{0.2pt}
\setlength{\headheight}{15pt}
\fancyhead[r]{\today}

\newcommand{\magnitude}[1]{\left\lvert #1 \right\rvert}

\begin{document}
\begin{enumerate}[label=\Alph*., leftmargin=*]
    \item Let $F:\mathbb{R}^{3}\rightarrow\mathbb{R}^2$ be defined by $F(x_1, x_2, x_3)=(\frac{x_1}{1+x_3},\frac{x_2}{1+x_3})$. Let $x$ by defined by $x(u,v)=(\frac{2u}{u^2+v^2+1},\frac{2v}{u^2+v^2+1},\frac{-(u^2+v^2-1)}{u^2+v^2+1})$.
          We compute that $F(x(u,v))=(u,v)$.
    \item \begin{enumerate}[label=.\Alph*]
        \item \textit{$x$ is regular when $\alpha$ is regular.}
        \begin{proof}
            We compute $\pderiv[x]{u}=(\pderiv[\alpha_1]{u},\pderiv[\alpha_2]{u},0)$ and $\pderiv[x]{v}=(0,0,1)$. For $\pderiv[x]{u}$ and $\pderiv[x]{v}$ to be linearly independent on all of the surface, we see that $(\alpha_1^{'}, \alpha_2^{'})\neq(0,0)$. So, it must be the case that $\alpha$ is regular.
        \end{proof}
        \item The image of $x$ is a surface when $\alpha$ behaves nicely. For example, if $\alpha$ self-intersects, the image of $x$ is not a surface.
    \end{enumerate}
    \item \begin{enumerate}[label=.\Alph*]
        \item Note that in general, $f(x)=x^{T}Mx=M_{11}{x_1}^2+M_{22}{x_2}^2+M_{33}{x_3}^2+2M_{12}x_1x_2+2M_{13}x_1x_3+2M_{23}x_2x_3$. When $M=\begin{bmatrix}
            0&0&0\\
            0&0&0\\
            0&0&1
        \end{bmatrix}$ then $f(x)=x_3^2=1$. This gives us an $x_1x_2$-plane at $x_3=\pm1$. When $M=\begin{bmatrix}
            0&0&0\\
            0&1&0\\
            0&0&1
        \end{bmatrix}$, $f(x)=1=x_2^2+x_3^2$. This creates a unbounded cylinder along the $x_1$ axis. When $M=\begin{bmatrix}
            1&0&0\\
            0&1&0\\
            0&0&1
        \end{bmatrix}$, then $f(x)=1=x_1^2+x_2^2+x_3^2$. This creates the unit sphere.
        \item \textit{$S=\{x|f(x)=x^{T}Mx=1$\} is non-empty when $M$ has at least one positive eigenvalue.} 
        \begin{proof}
            Let $\lambda$ be a positive eigenvalue of $M$, i.e. there is eigenvector $v$ that satisfies $Mv=\lambda v$. We can then consider $f(cv)$ with $c\in \mathbb{R}$. We compute\begin{equation*}
                f(cv)={(cv)}^{T}M(cv)=cv^{T}c\lambda v=c^2\magnitude{v}^2\lambda.
            \end{equation*} We want $c^2\magnitude{v}^2\lambda = 1$, so solving, we get that $c=\frac{1}{\sqrt{\magnitude{v}^2\lambda}}$. We know this $c$ exists because $\lambda$ is positive.
            So, when $x=\frac{v}{\sqrt{\magnitude{v}^2\lambda}}$ where $\lambda, v$ are a positive eigenvalue of $M$ and its corresponding eigenvector, $f(x)=1$, and therefore $S$ is non-empty.
        \end{proof}
    \end{enumerate}
\end{enumerate}
\end{document}